\section{Practices for Every Study}
\label{sec:crosscutting}

Ten practices apply whatever role the model plays. They are stated as
imperatives for brevity; the tier scheme of Section~\ref{sec:tiers} governs
how far each should be taken.

\subsection{Say What the Model Is Doing and What Is Being Claimed}

\begin{practice}{bp:declare}{Declare roles, levels, and claims}
State every role each model plays, in the robot and in the analysis; state at
which levels evidence was gathered; and state which claims the evidence is
offered for.
\end{practice}

The declaration costs a paragraph and removes the most common obstacle to
assessing an LLM-based HRI paper, which is not knowing what the model did.
Phrases such as ``the robot was powered by GPT-4'' conceal whether the model
generated every utterance verbatim, selected among scripted utterances,
paraphrased a script, produced labels that a rule-based dialogue manager
consumed, or drafted text that an experimenter approved. These are different
roles with different risks. The declaration should include roles in the
research pipeline, and it should say which outputs reached participants
without human review.

\subsection{Report the Configuration as the Stimulus}

\begin{practice}{bp:config}{Report the full configuration}
Report the model snapshot, provider, access dates, sampling parameters,
complete prompts, injected context, tools, memory, guardrails, and the speech
and perception components around the model, in enough detail for another
laboratory to reconstruct the stimulus.
\end{practice}

In an experiment the robot's behaviour is the stimulus, and with a generative
model the stimulus is produced by the conjunction of model weights, prompt,
context, and sampling parameters. Omitting any of them is equivalent to not
describing the stimulus. Three facts make this more than a formality. Hosted
models change: the measured behaviour of the same named commercial model on
the same tasks shifted substantially within months \cite{chen2024how}, so a
product name without a dated snapshot identifies a moving target. Hosted
models are withdrawn, which makes exact replication impossible and is a
principal argument for justifying any reliance on proprietary models in
research \cite{palmer2024using,ollion2024dangers,rogers2023closed}. And
outputs are sensitive to prompt details that authors regard as incidental:
changes of formatting alone, with meaning held constant, moved one model's
accuracy by as much as 76 points \cite{sclar2024quantifying}.

For robots, the configuration extends beyond the model. The system prompt
usually carries the robot's persona, its stated capabilities, and its
behavioural boundaries, and is therefore part of the design under study.
Expert consultation at RO-MAN 2025 produced a template with eight functional
components for specifying, bounding, and adapting robot behaviour through
prompts, proposed explicitly as a reporting aid as well as a design aid
\cite{ashok2026why}; we recommend its use. The pipeline around the model
(speech recogniser, voice-activity detection, text-to-speech, the mapping from
model output to gesture, any content filter and what the robot does when the
filter fires) shapes what the participant experiences as much as the model
does and should be reported with it. For action models, the checkpoint,
fine-tuning data, embodiment, control frequency, and the safety layer between
policy and actuators are the corresponding items. Experiment platforms that
store prompts and configurations alongside interaction logs make such
reporting a by-product of running the study \cite{laban2024lexi}.

\subsection{Measure Consistency Before Measuring Anything Else}

\begin{practice}{bp:consistency}{Measure run-to-run and prompt-to-prompt consistency}
Run the model more than once on the same inputs, and with reasonable variants
of the prompt, and report how much its outputs or scores vary.
\end{practice}

Psychometrics holds that reliability bounds validity: an instrument that does
not agree with itself cannot agree with anything else. A human coder is asked
to demonstrate intra-rater reliability by coding the same material twice. The
same should be asked of a model, and it frequently is not. Generative models
are stochastic by default, and configuring a hosted model to be
deterministic does not make it so: across ten runs of five such models,
accuracy varied by up to 15\% \cite{atil2025nondeterminism}.
Repeated classification of identical texts by the same model with the same
prompt can fall below conventional reliability thresholds
\cite{reiss2023testing}, model judges asked to rate the same output
repeatedly have shown low intra-rater reliability, in the worst case close to
arbitrary \cite{haldar2025rating,schroeder2024can}, and conclusions drawn
from a single prompt often do not survive paraphrase of the instruction: in
an evaluation of twenty models on thirty-nine tasks, paraphrasing the
instruction changed both absolute and relative performance
\cite{mizrahi2024state}.

In practice this means, for a model that labels or scores (\role{1},
\role{7}, \role{8}), running each item several times and reporting the
agreement among runs with the coefficient that will later be used for
agreement with people, so that the two are comparable. As a working default
we suggest five runs, increased when agreement among them is below the level
required of human coders. For a model that generates (\role{4}, \role{5}),
consistency concerns properties and not surface form: across repeated
generations from the same context, how often does the response satisfy each
criterion of the rubric, and how often does it contain a serious error? For a
model that decides or acts (\role{3}, \role{6}), it is the repeatability of
outcomes from nominally identical initial conditions. Agent benchmarks have
begun to report this as a quantity in its own right, the probability that a
task is completed on all of $k$ independent attempts, and find it far below
single-attempt success \cite{yao2025taubench}; gains in capability across
model generations have brought only small gains in reliability so defined
\cite{rabanser2026towards}. For a robot that someone will rely on daily,
succeeding every time is the property of interest. Where the study's conclusion depends on a
prompt, at least a small set of meaning-preserving variants should be tested
and the spread reported. Section~\ref{sec:gtheory} describes a design that
estimates these sources of variance jointly.

\subsection{Build a Human Reference and Show That It Is Reliable}

\begin{practice}{bp:human}{Establish a human reference with documented reliability}
Have at least two suitably qualified people independently rate or code the
material against a written rubric, report their agreement with an appropriate
coefficient and its confidence interval, and treat that agreement as the
ceiling against which the model is read.
\end{practice}

Human judgement is the reference for almost everything HRI wishes to know
about a model's outputs, and it is only as good as its reliability. Reviews
of human evaluation in natural language generation document how often this
is neglected: studies omit the number of raters, their expertise, the
question they were asked, and any measure of agreement
\cite{vanderlee2021human}. An annotation of 165 NLG papers found more than
200 different terms for the quality aspects being evaluated
\cite{howcroft2020twenty}, which later work has tried to consolidate into a
standard set of named and defined criteria \cite{belz2025standard}; a review
of 135 papers found that most reported agreement fell short of conventional
thresholds \cite{amidei2019agreement}; and when published human evaluations
were selected for reproduction, only 13\% provided enough information to
attempt it, and flaws were found in nearly all that were attempted
\cite{belz2023missing,thomson2024common}. Neighbouring fields are little better: in a sample
of qualitative papers from CSCW and HCI venues, inter-rater reliability was
reported in about one in nine \cite{mcdonald2019reliability}. HRI, which has
a long practice of behavioural coding, is well placed to do better, but the
practice needs adapting in four respects.

\paragraph{Choose the coefficient to fit the data.} Table~\ref{tab:irr}
summarises the choice. Percentage agreement should not be reported alone
because it does not correct for chance \cite{cohen1960coefficient,
hallgren2012computing}. Krippendorff's $\alpha$ handles any number of raters,
missing ratings, and nominal to ratio data, and is the most generally
applicable single choice \cite{hayes2007answering,krippendorff2019content}.
For continuous ratings the intraclass correlation is appropriate, provided
the model (one-way or two-way, random or mixed), the type (consistency or
absolute agreement), and the unit (single or average rating) are stated,
since these change the value and its meaning
\cite{shrout1979intraclass,koo2016guideline}.

\begin{table}[t]
\caption{Choosing an agreement statistic for a human reference, or for
comparing a model with one. In every case report the number of raters and
items, the prevalence of each category, and a confidence interval.}
\label{tab:irr}
\small
\begin{tabularx}{\textwidth}{@{}L{3.6cm}L{4.4cm}Y@{}}
\toprule
Data & Statistic & Notes \\
\midrule
Nominal labels, two raters & Cohen's $\kappa$
\cite{cohen1960coefficient} & Sensitive to prevalence and to differences in
raters' marginals; add per-category agreement \\
Nominal labels, more than two raters or missing ratings & Krippendorff's
$\alpha$ \cite{hayes2007answering}; Fleiss' $\kappa$
\cite{fleiss1971measuring} & $\alpha$ accommodates incomplete designs \\
Ordinal ratings (e.g.\ rubric scores) & Weighted $\kappa$
\cite{cohen1968weighted}; ordinal $\alpha$ & State the weighting scheme \\
Continuous or scale ratings & Intraclass correlation
\cite{shrout1979intraclass,koo2016guideline} & State model, type, and unit;
use absolute agreement when the values will be used as scores \\
Rare categories (e.g.\ safety violations) & $\kappa$ or $\alpha$ with
prevalence, positive and negative agreement; Gwet's AC1
\cite{gwet2008computing} & High raw agreement can coexist with low $\kappa$
\cite{feinstein1990high,byrt1993bias} \\
Two measurement methods on a continuous quantity & Bland--Altman limits of
agreement \cite{bland1986statistical}; concordance correlation
\cite{lin1989concordance} & Correlation alone ignores systematic offset \\
Interpretive qualitative analysis & Agreement statistics may be
inappropriate \cite{mcdonald2019reliability,braun2019reflecting} & Report the
analytic approach and how rigour was pursued instead \\
\bottomrule
\end{tabularx}
\end{table}

\paragraph{Interpret the value in context.} The verbal bands commonly
attached to $\kappa$ (``moderate'', ``substantial'') were offered by their
authors as arbitrary conveniences \cite{landis1977measurement}, and
Krippendorff's often-cited guidance to rely on variables with
$\alpha\geq.800$ and to draw only tentative conclusions from those between
.667 and .800 was offered for content analysis in general
\cite{krippendorff2019content}. Neither is a pass mark, and for subjective
judgements a fixed threshold is hard to defend at all; an alternative is to
benchmark agreement against a replication of the human rating itself
\cite{wong2021crossreplication}. What matters is the
consequence of the residual disagreement for the use to which the ratings are
put, and for that the confidence interval is as important as the point
estimate. Table~\ref{tab:kappa} in the Appendix shows that with a hundred
double-coded items, a common size for a validation subset in HRI, an observed
$\kappa$ of .60 is compatible with true values from roughly .44 to .74 when
categories are balanced, and with a much wider range when one category is
rare.

\paragraph{Watch prevalence.} The categories of greatest interest in HRI are
often rare: the unsafe utterance, the distressed child, the hazardous motion.
When a category is rare, raters who agree on nearly every item can have a low
$\kappa$, and, conversely, high raw agreement can conceal an inability to
agree on the rare cases at all \cite{feinstein1990high,byrt1993bias}. A
robot-led assessment of children's wellbeing with a vision-language model
illustrates the point from the model's side: the model agreed with a
psychologist on about 82\% of the items where no concern was coded, but on
only about 13\% of the items where a concern was coded, for a Cohen's
$\kappa$ of .22 \cite{abbasi2025robotled}. Overall agreement would have
suggested a serviceable instrument. Per-category agreement showed one that
missed most of what it was deployed to find.

\paragraph{Decide what disagreement means.} For some constructs disagreement
between raters is noise to be reduced by training and a sharper codebook. For
others, including many that HRI cares about (was that response empathic? was
that interruption rude? was that gesture appropriate?), disagreement reflects
real differences between people and cultures, and collapsing it into a
majority label discards information and installs one group's view as ground
truth \cite{aroyo2015truth,plank2022problem,uma2021learning}. For such
constructs we recommend reporting the distribution of ratings, recruiting
raters from the populations the robot is meant to serve, and evaluating
models against the distribution (Section~\ref{sec:perspectivist}).

Two further points concern who the raters are. Expertise should match the
construct: the clinical appropriateness of a wellbeing robot's response is a
matter for clinically trained raters, and lay raters have been shown to be
unreliable judges even of whether text was written by a person
\cite{clark2021all} and of the quality of open-ended generated text
\cite{karpinska2021perils}. And raters should be blind to the source of what
they rate (model or human, and which condition), since knowledge of the
source shapes judgement.

\subsection{Validate Automated Evaluators Before Relying on Them}

\begin{practice}{bp:validate}{Validate every automated evaluator against the human reference}
Before a metric, a model judge, or a model annotator is used to support a
claim, demonstrate on held-out, human-rated material from the study's own
domain that it agrees with people about as well as people agree with each
other.
\end{practice}

Automatic evaluation is attractive in HRI for the same reason as elsewhere:
human rating is slow and expensive. Its history in dialogue research is a
caution. Word-overlap metrics imported from machine translation were found to
correlate weakly or not at all with human judgements of dialogue responses
\cite{liu2016how}, yet were reported for years. Model-based judges are far
more capable, and on some tasks reach agreement with human preferences
comparable to agreement between humans \cite{zheng2023judging,chiang2023can}.
But the evidence for them is uneven across tasks: a study spanning twenty NLP
datasets found large variation in how well model judgements tracked human
ones and concluded that judges need to be validated for each task before use
\cite{bavaresco2025llms}. Even where humans agree strongly with one
another, only the best of thirteen judge models aligned reasonably with them
and none reached the level of agreement among humans, and judges with high
percentage agreement still assigned very different scores
\cite{thakur2025judging}. Without a correct reference answer, judges agreed
with experts mainly on the questions they could answer themselves
\cite{krumdick2026no}. Judges also have systematic biases, including
sensitivity to the order in which candidates are presented
\cite{wang2024large}, a preference for text produced by the judging model
itself \cite{panickssery2024llm}, and a range of other cognitive-bias
analogues \cite{koo2024benchmarking}.

Validation on HRI material has only begun, and what there is counsels
caution. When language models were asked to reproduce participants'
responses in three published HRI user studies, the strongest model
correlated highly with people in two of them and no model did in the third;
vision-capable models failed on the video stimuli; and models rated robot
behaviour more favourably than people had \cite{wachowiak2024are}. When
models completed standard HRI questionnaires for a set of recorded
interactions, they agreed moderately to strongly with human observers on
engagement and were highly consistent from run to run, yet their ratings of
comfort and strangeness were \emph{inversely} related to those of people
\cite{lockwood2026when}. The second result deserves emphasis because it
separates two things that are easily conflated: the models were reliable,
and on that dimension they were wrong. The constructs that HRI most often
asks a judge to score (social appropriateness in a physical context,
fidelity to a therapeutic technique, consistency between an utterance and
what the robot's body is doing) and the inputs it most often supplies
(transcripts of recognised speech with errors, video of a robot and a
person) remain largely unvalidated. A judge's validity is specific to the construct, the rubric, the
prompt, and the kind of material. It has to be established locally, on the
study's own data, with the study's own rubric, and the role-specific
practices for Annotators and Judges (\bpr{bp:coder}, \bpr{bp:judge}) say how.

\subsection{Quantify Uncertainty and Plan for Precision}

\begin{practice}{bp:uncertainty}{Report uncertainty and size evaluations deliberately}
Attach an interval to every component-level estimate, choose the number of
items, trials, raters, and runs in advance on the basis of the precision
required, and analyse nested data as nested.
\end{practice}

Component-level results in HRI papers are frequently single numbers from
small samples. The statistical literatures of both NLP and robot learning have
argued at length that this is inadequate: significance testing and power
analysis are underused in NLP \cite{dror2018hitchhikers,card2020with}, model
evaluations are experiments whose results need standard errors
\cite{miller2024adding}, error bars based on the central limit theorem
understate uncertainty for evaluation sets of fewer than a few hundred items
\cite{bowyer2025position}, and reports of robot policy performance from a
handful of trials do not support the comparisons drawn from them
\cite{kressgazit2024robot,agarwal2021deep}. Table~\ref{tab:cp} in the Appendix
makes the last point concrete.

Two features of HRI data need particular care. The data are nested: turns
within sessions within participants, trials within objects within scenes.
Treating five hundred robot utterances from twenty participants as five
hundred independent observations overstates precision; mixed-effects models
or cluster bootstraps are the appropriate tools, as they are for the
participant-level outcomes of the same studies. And evaluation sets are
typically built after the fact from whatever the study produced. Deciding in
advance how many items will be double-coded, how they will be sampled
(stratified by session, participant, and condition, with over-sampling of rare
categories), and what interval width is acceptable turns validation from a
ritual into a measurement.

\subsection{Disaggregate}

\begin{practice}{bp:subgroup}{Report performance for the groups of people the robot serves}
Break component-level performance down by participant characteristics
relevant to the application, and test whether error rates differ.
\end{practice}

Aggregate performance can hide failure concentrated in a subgroup, and
embodiment gives such failure direct consequences. Robots driven by
foundation models have been shown to act on social stereotypes
\cite{hundt2022robots}, and LLMs prompted in HRI scenarios produce
discriminatory outputs across protected characteristics and accept harmful
instructions \cite{hundt2025llm}. In the child wellbeing assessment discussed
above, the vision-language model's false-positive rate was significantly
higher for girls than for boys, so that it failed the equalised-odds criterion
even when gender was not mentioned in the prompt \cite{abbasi2025robotled}.
Speech recognition, which sits upstream of most spoken HRI, has well-documented
disparities by speaker group \cite{koenecke2020racial}, and an error there
becomes an error in everything the language model does next. HRI samples are
small and not diverse \cite{seaborn2023not}, which limits what
disaggregation can show; but reporting the breakdown, with intervals, allows
later synthesis, and its absence should be stated as a limitation and not
passed over.

\subsection{Characterise Failures, Not Only Averages}

\begin{practice}{bp:failures}{Report an error taxonomy, severity, and the worst cases}
Classify the model's errors by type and by severity for the person, report the
frequency of each, and describe the most serious individually.
\end{practice}

A mean rubric score of 4.3 out of 5 is uninformative about whether the robot
once told a participant something dangerous. HRI has taxonomies of robot
failure and of social error that can be extended to generative components
\cite{honig2018understanding,tian2021taxonomy}, and the NLP literature
supplies taxonomies of hallucination \cite{ji2023survey}. Failure modes that
studies of LLM-driven robots have documented, and that a taxonomy for HRI
should therefore include, are fabricated facts and memories, claims to
capabilities or experiences the robot does not have, loss of persona or of
conversational context, responses that are too long or too late for spoken
interaction, misrecognised speech treated as intended, and inappropriate
advice \cite{irfan2025between,ashok2026why}. For action models the
corresponding list includes hesitation, motion that is successful but
illegible or alarming, and task success achieved in a way that would be
unacceptable around a person \cite{guan2024task}. Severity should be rated
separately from frequency and with reference to the population: the same
error has different weight for a student volunteer and for a person with
dementia.

\subsection{Connect the Levels}

\begin{practice}{bp:levels}{Audit what participants actually received, and measure the system around the model}
Log every model output delivered in the study and evaluate those outputs, or a
defined sample of them, at component level; report system-level properties
that mediate between model and person; and verify manipulations in the
delivered outputs.
\end{practice}

This practice answers the two inferential leaps of
Section~\ref{sec:levels}. Against the downward leap, it requires that
outputs produced \emph{in situ} be audited, since outputs produced at the
desk before the study are not the outputs participants heard. A generative
model makes every participant's interaction different; the audit is the only
means of knowing what the ``treatment'' was. It also supplies the
manipulation check that experimental method requires
\cite{hoffman2021primer}: if two conditions differ in a prompt intended to
make the robot more empathic, the audit should show that raters blind to
condition find its utterances more empathic.

Against the upward leap, it requires measurement of what lies between model
and person. Response latency is the clearest case. Gaps between turns in
human conversation are typically a few hundred milliseconds across languages
\cite{stivers2009universals}, and timing is central to how spoken interaction
with machines is experienced \cite{skantze2021turn}; a pipeline of speech
recognition, a hosted model, and speech synthesis commonly takes several
times longer, and the delay, the filler that masks it, and the interruptions
it causes are part of the stimulus. The distribution of latency, and not only
its mean, should be reported. So should the word error rate of the speech
recogniser in the study's acoustic conditions and with the study's
population, since recognition errors are a leading cause of breakdown in
LLM-driven conversation with older adults \cite{irfan2025between}.

Where a human supervises the robot, their interventions are data. Riek's
guidelines for Wizard-of-Oz studies asked authors to report what the wizard
could do, what they did, and how they were trained \cite{riek2012wizard}.
Supervised autonomous systems need the equivalent: under what conditions the
supervisor could veto, edit, or replace an output, how often they did, and
why. The intervention rate is among the most informative system-level
measures of a generative component's readiness, and it is almost free to
collect.

\subsection{Preserve, Share, and Pre-Specify}

\begin{practice}{bp:open}{Archive the materials and fix the evaluation plan in advance}
Preserve prompts, configurations, interaction logs, rater data, and judge
outputs; share what consent and privacy allow; and specify prompts, rubrics,
judges, validation subsets, and thresholds before the evaluation is run.
\end{practice}

Generative models multiply the analytic choices available to a researcher.
The prompt can be reworded, the model swapped, the temperature changed, the
rubric adjusted, the judge replaced, each change producing a somewhat
different result. The lesson of the replication crisis in psychology is that
undisclosed flexibility of this kind reliably produces false positives
\cite{simmons2011falsepositive}. The same has now been shown for model
configurations. In a replication of 37 annotation tasks from 21 social
science studies, using 18 models and some 13 million labels, the choice
among defensible configurations led to incorrect conclusions for roughly a
third of hypotheses even with state-of-the-art models, and a few
paraphrases of the prompt were enough to make almost any result significant
\cite{baumann2025large}. Re-running one published use of simulated
respondents under 66 defensible configurations moved its correlation with
human data from .23 to .84 \cite{cummins2026threat}. Reporting results
across the set of reasonable prompts, models, and decoding settings, as a
multiverse or specification curve
\cite{steegen2016increasing,simonsohn2020specification}, is the transparent
alternative to reporting the configuration that happened to work.
Pre-registration, the remedy psychology adopted \cite{nosek2018preregistration},
transfers directly: the prompt, the model, the rubric, the validation sample,
and the agreement threshold can be fixed and time-stamped before the model is
run on the study data. Prompt development is legitimate, but it should be
done on separate material, and the paper should say so.

Preservation matters more than usual because the model itself cannot be
archived when it is hosted. The interaction logs then become the only
permanent record of what the system did, and a set of archived contexts is
what makes later comparison with a successor model possible
(Section~\ref{sec:replay}). HRI's reproducibility agenda already calls for
sharing code, materials, and data \cite{gunes2022reproducibility}; for
studies with generative models the minimum addition is the complete prompt
set and, where consent permits, de-identified logs. Sending participants'
speech, images, or disclosures to a third-party model provider is a
data-protection matter that belongs in ethics review and in the consent form,
and is a reason, alongside reproducibility, to consider locally run
open-weight models where they are adequate to the role, as was done for the
assessment of children's narratives cited above \cite{abbasi2025robotled}.
